\lyseis{Β1.12}{Επίκεντρη γωνία, τόξο και μέτρηση τόξου}
\sol{1} Μικρό τόξο $70^\circ$, μεγάλο $360^\circ-70^\circ=290^\circ$.
\sol{2} $360^\circ:5=72^\circ$, $360^\circ:9=40^\circ$, $360^\circ:12=30^\circ$.
\sol{3} $90^\circ$, $45^\circ$, $270^\circ$, $300^\circ$.
\sol{4} $\frac{60}{360}=\frac{1}{6}$ του κύκλου.
\sol{5} Κάθε ημικύκλιο είναι $180^\circ$.
\sol{6} Τα τόξα είναι $100^\circ$, $120^\circ$, $140^\circ$ και το άθροισμά τους $360^\circ$ (όλος ο κύκλος).
\sol{7} α) Σωστό. β) Σωστό. γ) Λάθος: εξαρτάται από την ακτίνα. δ) Σωστό.
\sol{8} $360^\circ:12=30^\circ$.
\sol{9} $20$ λεπτά: $120^\circ$. $45$ λεπτά: $270^\circ$. Μία ώρα: $360^\circ$.
\sol{10} $360:45=8$ φέτες. Τρεις φέτες είναι τα $\frac{3}{8}$ της πίτας.
\sol{11} $\frac{75}{360}=\frac{5}{24}$ του κύκλου.
\sol{12} $4$ τόξα.
\sol{13} Τα μέτρα είναι ίσα (ίδια επίκεντρη γωνία). Τα μήκη δεν είναι ίσα: το τόξο του μεγαλύτερου κύκλου είναι μακρύτερο.
\sol{14} Ο $360$ έχει πολλούς διαιρέτες και ταιριάζει με το $60$ των Βαβυλωνίων. Αν επιλέγαμε $100$ θα άλλαζαν οι αριθμοί αλλά όχι η γεωμετρία (αυτή η μονάδα λέγεται βαθμός grad).
\sol{15} $360^\circ-120^\circ=240^\circ$.
\sol{16} Διπλασιάζεται και το τόξο.
\sol{17} $360:6=60$ μοίρες.
\sol{18} Ελεύθερη απάντηση.
